\chapter{Critical Evaluation}
\label{chap:evaluation}

% ============================================================
\section{Overview: Reproduction of Stage (a) and Stage (b)}
\label{sec:repro_overview}

\noindent
This chapter provides a critical evaluation of the reproduction results
for Stage~(a) (standard supervised learning) and Stage~(b)
(semi-supervised and multi-task learning) based on the framework proposed
by \cite{hoyer2023ijcv_sde}, followed by an evaluation of the proposed
Cross-Task Consistency Loss ($\mathcal{L}_\text{ct}$). The primary
objective is to verify whether the reported performance trends are
consistent under a constrained academic computing environment and whether
$\mathcal{L}_\text{ct}$ yields reliable improvements over the reproduced
multi-task baseline. All experiments were conducted on the Cityscapes
dataset, with results reported as mean mIoU~(\%) $\pm$ standard deviation
across three independent seeds (7, 25, 42). To ensure transparency, each
reproduction entry is presented as ours / (paper), followed by the
absolute discrepancy $\Delta$ (ours $-$ paper).

% ============================================================
\section{Reproduction Protocol}
\label{sec:repro_protocol}

\noindent
To avoid redundancy, this evaluation follows the design and implementation
described in Chapter~\ref{chap:execution}. Stage~(a) and transfer
experiments use a ResNet-101 encoder with an ASPP-based decoder;
Stage~(b) multi-task and cross-task loss experiments additionally employ
the PAD decoder (Section~\ref{sec:bridge}). We used the author-provided
fixed subset indices for the 100, 372, and 744 label regimes. All models
were trained for 40,000 iterations using SGD with polynomial
learning-rate decay, and results are reported as mean mIoU~(\%) $\pm$
standard deviation over three seeds. For Stage~(b) transfer-related
experiments (Table~5), we use the author-released SDE checkpoints
(fd0/fd2) to maintain parity with the official pipeline.

Baseline values are reported per table under the corresponding experiment
configuration. When baseline values differ across tables, this reflects
differences in the surrounding training pipeline rather than implying
direct numerical comparability across tables.

% ============================================================
\section{Stage (a): Data Selection Reproduction (Table 1)}
\label{sec:stage_a_table1}

\noindent
\textbf{Objective.} Stage~(a) assesses the effectiveness of various data
selection strategies under limited label budgets.
Table~\ref{tab:table1_repro} compares the random baseline against
entropy-based selection and the SDE-guided variants (US, DS, and DS+US).

\begin{table}[H]
\centering
\caption{Reproduction of Table~1 (Data Selection). mIoU~(\%) as
mean$\pm$std (3 seeds). Each cell shows: ours / {\scriptsize (paper)},
$\Delta$.}
\label{tab:table1_repro}
\renewcommand{\arraystretch}{1.6}
\small
\begin{tabularx}{\textwidth}{lXXX}
\toprule
\textbf{Method} &
\makecell[c]{\textbf{1/30} \\ \textbf{(100 labels)}} &
\makecell[c]{\textbf{1/8} \\ \textbf{(372 labels)}} &
\makecell[c]{\textbf{1/4} \\ \textbf{(744 labels)}} \\
\midrule
Random &
\makecell[c]{$47.57 \pm 0.96$ \\ {\scriptsize ($48.75 \pm 1.61$), $\Delta=-1.18$}} &
\makecell[c]{$58.63 \pm 0.67$ \\ {\scriptsize ($59.14 \pm 1.02$), $\Delta=-0.51$}} &
\makecell[c]{$63.22 \pm 0.30$ \\ {\scriptsize ($63.46 \pm 0.38$), $\Delta=-0.24$}} \\
Entropy &
\makecell[c]{$53.76 \pm 0.44$ \\ {\scriptsize ($53.63 \pm 0.77$), $\Delta=+0.13$}} &
\makecell[c]{$63.61 \pm 0.78$ \\ {\scriptsize ($63.51 \pm 0.68$), $\Delta=+0.10$}} &
\makecell[c]{$65.57 \pm 0.62$ \\ {\scriptsize ($66.18 \pm 0.50$), $\Delta=-0.61$}} \\
Ours (US) &
\makecell[c]{$50.83 \pm 0.91$ \\ {\scriptsize ($51.75 \pm 1.12$), $\Delta=-0.92$}} &
\makecell[c]{$63.07 \pm 0.11$ \\ {\scriptsize ($62.77 \pm 0.46$), $\Delta=+0.30$}} &
\makecell[c]{$66.84 \pm 0.34$ \\ {\scriptsize ($66.76 \pm 0.45$), $\Delta=+0.08$}} \\
Ours (DS) &
\makecell[c]{$52.67 \pm 0.14$ \\ {\scriptsize ($53.00 \pm 0.51$), $\Delta=-0.33$}} &
\makecell[c]{$62.86 \pm 0.46$ \\ {\scriptsize ($63.23 \pm 0.69$), $\Delta=-0.37$}} &
\makecell[c]{$66.74 \pm 0.15$ \\ {\scriptsize ($66.37 \pm 0.20$), $\Delta=+0.37$}} \\
Ours (DS+US) &
\makecell[c]{$53.63 \pm 0.19$ \\ {\scriptsize ($54.37 \pm 0.36$), $\Delta=-0.74$}} &
\makecell[c]{$63.38 \pm 0.92$ \\ {\scriptsize ($64.25 \pm 0.18$), $\Delta=-0.87$}} &
\makecell[c]{$66.35 \pm 0.62$ \\ {\scriptsize ($66.94 \pm 0.59$), $\Delta=-0.59$}} \\
\bottomrule
\end{tabularx}
\end{table}

\noindent
\textbf{Observation and Analysis.} Our reproduction confirms the overall
trend that informed selection strategies consistently outperform Random
across all label budgets. Relative to Random, entropy and the
depth-guided variants yield substantial gains under the 100 and 372
label regimes, while the margin narrows at 744 labels, suggesting
diminishing returns as supervision increases. Although the paper reports
DS+US as the most consistently effective strategy, our results indicate
that Entropy can achieve comparable performance in the lowest-budget
settings. The remaining absolute gaps to the paper (typically within
$\sim$1 mIoU) may be partly attributable to single-GPU constraints
(e.g., small batch size without SyncBN), but could also reflect
run-to-run variance.

\clearpage
% ============================================================
\section{Stage (b): Data Mixing Strategy Reproduction (Table 3)}
\label{sec:stage_ab_table3}

\noindent
\textbf{Objective.} This section evaluates Stage~(b) mixing strategies
applied to the subsets defined in Stage~(a).
Table~\ref{tab:table3_repro} compares the baseline against
pseudo-labelling, ClassMix, and the proposed geometry-aware
DepthMix~(DX).

\begin{table}[H]
\centering
\caption{Reproduction of Table~3 (Mixing Strategies). mIoU~(\%)
mean$\pm$std (3 seeds). $\Delta$ reflects deviation from the paper.}
\label{tab:table3_repro}
\renewcommand{\arraystretch}{1.2}
\begin{tabularx}{\textwidth}{lXX}
\toprule
\textbf{Mixing Strategy} &
\textbf{372 Labels (1/8)} &
\textbf{2975 Labels (Full)} \\
\midrule
Baseline &
$58.25 \pm 0.43$ {\scriptsize ($59.14 \pm 1.02$), $\Delta=-0.89$} &
$67.30 \pm 0.53$ {\scriptsize ($67.77 \pm 0.13$), $\Delta=-0.47$} \\
Pseudo-Labels &
$60.05 \pm 0.92$ {\scriptsize ($62.39 \pm 0.86$), $\Delta=-2.34$} &
-- \\
ClassMix &
$63.46 \pm 1.20$ {\scriptsize ($63.16 \pm 0.89$), $\Delta=+0.30$} &
$69.23 \pm 0.51$ {\scriptsize ($69.60 \pm 0.32$), $\Delta=-0.37$} \\
DepthMix &
$64.57 \pm 1.12$ {\scriptsize ($64.14 \pm 1.34$), $\Delta=+0.43$} &
$70.11 \pm 0.28$ {\scriptsize ($69.83 \pm 0.36$), $\Delta=+0.28$} \\
\bottomrule
\end{tabularx}
\end{table}

\noindent
\textbf{Observation and Analysis.} The performance hierarchy
(DepthMix $>$ ClassMix $>$ Baseline) is effectively reproduced. The
superiority of DepthMix confirms that incorporating geometric scene
structure reduces boundary artefacts typical of ClassMix. The discrepancy
in pseudo-label results ($\Delta = -2.34$) underlines the sensitivity of
self-training to confidence thresholds and batch size, which were
restricted by hardware constraints.

% ============================================================
\section{Stage (b): SDE Feature Transfer Reproduction (Table 5)}
\label{sec:stage_ab_table5}

\noindent
\textbf{Objective.} We evaluate the impact of transferring SDE-pretrained
features to the segmentation task, specifically investigating the role of
the feature-distance loss ($\mathcal{L}_F$) and multi-task learning.

\begin{table}[H]
\centering
\caption{Reproduction of Table~5 (SDE Feature Transfer). mIoU~(\%)
mean$\pm$std (3 seeds).}
\label{tab:table5_repro}
\renewcommand{\arraystretch}{1.2}
\begin{tabularx}{\textwidth}{lXX}
\toprule
\textbf{Method} &
\textbf{372 Labels (1/8)} &
\textbf{2975 Labels (Full)} \\
\midrule
Baseline &
$58.25 \pm 0.43$ {\scriptsize ($59.14 \pm 1.02$), $\Delta=-0.89$} &
$67.30 \pm 0.53$ {\scriptsize ($67.77 \pm 0.13$), $\Delta=-0.47$} \\
Transfer (no $\mathcal{L}_F$) &
$59.73 \pm 0.57$ {\scriptsize ($60.46 \pm 0.64$), $\Delta=-0.73$} &
$68.74 \pm 0.67$ {\scriptsize ($69.00 \pm 0.70$), $\Delta=-0.26$} \\
Transfer ($\mathcal{L}_F=\checkmark$) &
$60.83 \pm 0.94$ {\scriptsize ($60.80 \pm 0.69$), $\Delta=+0.03$} &
$68.99 \pm 0.27$ {\scriptsize ($69.47 \pm 0.38$), $\Delta=-0.48$} \\
Multi-Task ($\mathcal{L}_F=\checkmark$) &
$60.75 \pm 0.69$ {\scriptsize ($61.25 \pm 0.55$), $\Delta=-0.50$} &
$69.40 \pm 0.29$ {\scriptsize ($69.76 \pm 0.39$), $\Delta=-0.36$} \\
\bottomrule
\end{tabularx}
\end{table}

\noindent
\textbf{Observation and Analysis.} Our results reinforce the core
scientific claims: SDE pretraining provides a superior initialisation,
and $\mathcal{L}_F$ successfully prevents catastrophic forgetting of
semantic features. The near-perfect alignment for Transfer
($\mathcal{L}_F=\checkmark$) at 372 labels ($\Delta = +0.03$) suggests
that learned geometric representations transfer robustly across training
environments. These results also motivate the design of
$\mathcal{L}_\text{ct}$: even in the MTL setting, depth decoder features
are not explicitly required to align with segmentation features, leaving
room for further cross-task regularisation.

% ============================================================
\section{Stage (b): Framework Component Ablation (Table 7)}
\label{sec:stage_b_table7}

\noindent
\textbf{Objective.} We reproduce the full framework component ablation
study to evaluate the synergistic effects of data selection~(S),
DepthMix~(DX), and multi-task learning~(MTL).
Table~\ref{tab:table7_repro} presents a direct comparison between our
results and the values reported in \cite{hoyer2023ijcv_sde}.

\begin{table}[H]
\centering
\caption{Reproduction of Table~7 (Framework Ablation). mIoU~(\%)
mean$\pm$std (3 seeds). Each cell shows: ours / {\scriptsize (paper)},
$\Delta$.}
\label{tab:table7_repro}
\renewcommand{\arraystretch}{2.0}
\small
\begin{tabularx}{\textwidth}{cccXX}
\toprule
\textbf{S} & \textbf{DX} & \textbf{MTL} &
\makecell[c]{\textbf{372 Labels (1/8)}} &
\makecell[c]{\textbf{2975 Labels (Full)}} \\
\midrule
-- & -- & -- &
\makecell[c]{$58.25 \pm 0.43$ \\ {\scriptsize ($59.14 \pm 1.02$), $\Delta=-0.89$}} &
\makecell[c]{$67.30 \pm 0.53$ \\ {\scriptsize ($67.77 \pm 0.13$), $\Delta=-0.47$}} \\
-- & -- & \checkmark &
\makecell[c]{$60.75 \pm 0.69$ \\ {\scriptsize ($61.25 \pm 0.55$), $\Delta=-0.50$}} &
\makecell[c]{$69.40 \pm 0.29$ \\ {\scriptsize ($69.76 \pm 0.39$), $\Delta=-0.36$}} \\
-- & \checkmark & -- &
\makecell[c]{$64.57 \pm 1.12$ \\ {\scriptsize ($64.14 \pm 1.34$), $\Delta=+0.43$}} &
\makecell[c]{$70.11 \pm 0.28$ \\ {\scriptsize ($69.83 \pm 0.36$), $\Delta=+0.28$}} \\
\checkmark & -- & -- &
\makecell[c]{$65.23 \pm 0.46$ \\ {\scriptsize ($64.25 \pm 0.18$), $\Delta=+0.98$}} &
\makecell[c]{--} \\
\checkmark & -- & \checkmark &
\makecell[c]{$66.81 \pm 0.28$ \\ {\scriptsize ($65.35 \pm 0.10$), $\Delta=+1.46$}} &
\makecell[c]{--} \\
\checkmark & \checkmark & -- &
\makecell[c]{$66.99 \pm 0.44$ \\ {\scriptsize ($66.48 \pm 0.27$), $\Delta=+0.51$}} &
\makecell[c]{--} \\
-- & \checkmark & \checkmark &
\makecell[c]{$65.11 \pm 0.51$ \\ {\scriptsize ($66.66 \pm 1.05$), $\Delta=-1.55$}} &
\makecell[c]{$71.28 \pm 0.30$ \\ {\scriptsize ($71.16 \pm 0.16$), $\Delta=+0.12$}} \\
\checkmark & \checkmark & \checkmark &
\makecell[c]{$67.66 \pm 0.47$ \\ {\scriptsize ($68.01 \pm 0.83$), $\Delta=-0.35$}} &
\makecell[c]{--} \\
\bottomrule
\end{tabularx}
\end{table}

\noindent
\textbf{Observation and Synergy Analysis.}
The full ablation confirms that S+DX+MTL achieves the highest performance
at the 372-label regime ($67.66$ mIoU), consistent with the paper's
findings. Our reproduced values are broadly well-aligned with the paper,
with most configurations falling within $\pm$1 mIoU.

A notable discrepancy is observed in the DX+MTL configuration for 372
labels, where our result ($65.11$) falls $1.55$ mIoU below the paper
($66.66$), despite exceeding the DX-only result ($64.57$). This
suggests that under our single-GPU setting, joint optimisation of depth
and segmentation is less effective when DepthMix augmentation is combined
with the auxiliary depth objective, compared to the multi-GPU setup
reported in the paper. This gap motivates the design of
$\mathcal{L}_\text{ct}$: rather than relying solely on the implicit
cross-task interaction in PAD, an explicit feature-alignment loss provides
a direct and controllable mechanism for transferring geometric knowledge
to the segmentation branch.

% ============================================================
\section{Cross-Task Consistency Loss: Calibration and Ablation}
\label{sec:ct-calibration-ablation}

Building on the reproduced PAD-MTL baseline, this section evaluates the
proposed Cross-Task Consistency Loss ($\mathcal{L}_\text{ct}$), which
aligns intermediate features from the segmentation and depth decoder
branches. We report results from two experimental phases: a single-seed
calibration sweep (Phase~I) over loss type, projection layout, warmup
schedule, and $\lambda_\text{ct}$; and a three-seed validation
(Phase~II) for the most promising configurations. All $\Delta$ values are
reported as absolute mIoU differences in percentage points~(pp).

% ============================================================
\subsection{Phase I: Single-seed Calibration (Seed~7)}
\label{subsec:ct-phase1-seed7}

To limit compute while exploring a large discrete grid, we fix random
seed~7 and measure best validation mIoU for every combination of:
loss type (cosine vs.\ normalised MSE), projection layout (single vs.\
dual learnable projections), warmup schedule ($w5k$: 5000
segmentation-only iterations before activating $\mathcal{L}_\text{ct}$;
$w0$: no warmup), and
$\lambda_\text{ct}\in\{0.25, 0.50, 0.75, 1.00\}$.

\paragraph{Baseline reference.}
The PAD-MTL baseline (Table~5, Multi-Task, 3~seeds) yields
seed-7 / 25 / 42 $=$ 59.84 / 60.92 / 61.50\,\%,
mean $= 60.75$\,\%, std $= 0.69$\,\%.
All Phase~I $\Delta$ values use the same-seed value
$\text{baseline}_{\text{S7}} = 59.84$\,\%.

\paragraph{Single-projection scan.}
Tables~\ref{tab:ct-p1-single-cos} and~\ref{tab:ct-p1-single-mse}
report the single-projection results. The single+MSE block was
independently re-run (Slurm jobs 16513723--16513730); all other blocks
follow the exp~219 sweep protocol.

\begin{table}[H]
\centering
\caption{Phase~I (seed~7): single projection + cosine.
$\Delta = \text{mIoU}(\%) - 59.84$\,pp.}
\label{tab:ct-p1-single-cos}
\renewcommand{\arraystretch}{1.2}
\begin{tabular}{llcc}
\toprule
Warmup & $\lambda_\text{ct}$ & Best mIoU (\%) & $\Delta$\,(pp) \\
\midrule
Baseline (seed~7) & -- & 59.84 & -- \\
\midrule
$w5k$ & 0.25 & 62.64 & $+2.80$ \\
$w5k$ & 0.50 & 61.69 & $+1.85$ \\
$w5k$ & 0.75 & 61.86 & $+2.02$ \\
$w5k$ & 1.00 & 61.40 & $+1.56$ \\
\midrule
$w0$  & 0.25 & 62.89 & $+3.05$ \\
$w0$  & 0.50 & 61.42 & $+1.58$ \\
$w0$  & 0.75 & 63.00 & $+3.16$ \\
$w0$  & 1.00 & 63.77 & $+3.93$ \\
\bottomrule
\end{tabular}
\end{table}

\begin{table}[H]
\centering
\caption{Phase~I (seed~7): single projection + MSE (rerun).
$\Delta = \text{mIoU}(\%) - 59.84$\,pp.}
\label{tab:ct-p1-single-mse}
\renewcommand{\arraystretch}{1.2}
\begin{tabular}{llcc}
\toprule
Warmup & $\lambda_\text{ct}$ & Best mIoU (\%) & $\Delta$\,(pp) \\
\midrule
Baseline (seed~7) & -- & 59.84 & -- \\
\midrule
$w5k$ & 0.25 & 62.58 & $+2.74$ \\
$w5k$ & 0.50 & 62.41 & $+2.57$ \\
$w5k$ & 0.75 & 62.54 & $+2.70$ \\
$w5k$ & 1.00 & 63.43 & $+3.59$ \\
\midrule
$w0$  & 0.25 & 61.97 & $+2.13$ \\
$w0$  & 0.50 & 61.68 & $+1.84$ \\
$w0$  & 0.75 & 61.46 & $+1.62$ \\
$w0$  & 1.00 & 62.16 & $+2.32$ \\
\bottomrule
\end{tabular}
\end{table}

\noindent
\textbf{Phase-I findings (single projection).}
\begin{itemize}
\item $\mathcal{L}_\text{ct}$ is consistently beneficial: all
  configurations improve over the seed-7 baseline under both loss types,
  with gains ranging from $+1.56$ to $+3.93$\,pp.
\item For cosine, $w0$ is clearly superior: the best cell is
  $w0$, $\lambda_\text{ct}=1.00$ ($63.77$\,\%, $+3.93$\,pp). With $w5k$,
  performance plateaus or degrades at larger $\lambda_\text{ct}$.
\item For MSE, $w5k$ is substantially better than $w0$ at
  $\lambda_\text{ct}=1.00$ ($63.43$\,\% vs.\ $62.16$\,\%, gap of $1.27$\,pp).
  A brief warmup stabilises the MSE gradient scale before the cross-task
  term activates.
\item These results narrow the candidate configurations for Phase~II:
  cosine with $w0$ and $\lambda_\text{ct}\approx1.0$; MSE with $w5k$
  and $\lambda_\text{ct}=1.0$.
\end{itemize}

% ============================================================
\subsection{Phase II: Three-seed Validation and Dual-Projection Comparison}
\label{subsec:ct-phase2-3seeds}

Using the candidates identified in Phase~I, we train with three
independent seeds (7, 25, 42) and additionally include the dual-projection
variants to assess the benefit of a shared versus independent projection
architecture.

\paragraph{Dual-projection reference (seed~7).}
Tables~\ref{tab:ct-p1-dual-cos} and~\ref{tab:ct-p1-dual-mse} provide
the dual-projection seed-7 scan from Phase~I for direct comparison
against the single-projection results above.

\begin{table}[H]
\centering
\caption{Phase~I (seed~7): dual projection + cosine.
$\Delta = \text{mIoU}(\%) - 59.84$\,pp.}
\label{tab:ct-p1-dual-cos}
\renewcommand{\arraystretch}{1.2}
\begin{tabular}{llcc}
\toprule
Warmup & $\lambda_\text{ct}$ & Best mIoU (\%) & $\Delta$\,(pp) \\
\midrule
Baseline (seed~7) & -- & 59.84 & -- \\
\midrule
$w5k$ & 0.25 & 62.80 & $+2.96$ \\
$w5k$ & 0.50 & 62.49 & $+2.65$ \\
$w5k$ & 0.75 & 62.67 & $+2.83$ \\
$w5k$ & 1.00 & 61.27 & $+1.43$ \\
\midrule
$w0$  & 0.25 & 61.76 & $+1.92$ \\
$w0$  & 0.50 & 63.66 & $+3.82$ \\
$w0$  & 0.75 & 62.45 & $+2.61$ \\
$w0$  & 1.00 & 63.32 & $+3.48$ \\
\bottomrule
\end{tabular}
\end{table}

\begin{table}[H]
\centering
\caption{Phase~I (seed~7): dual projection + MSE.
$\Delta = \text{mIoU}(\%) - 59.84$\,pp.}
\label{tab:ct-p1-dual-mse}
\renewcommand{\arraystretch}{1.2}
\begin{tabular}{llcc}
\toprule
Warmup & $\lambda_\text{ct}$ & Best mIoU (\%) & $\Delta$\,(pp) \\
\midrule
Baseline (seed~7) & -- & 59.84 & -- \\
\midrule
$w5k$ & 0.25 & 61.83 & $+1.99$ \\
$w5k$ & 0.50 & 62.67 & $+2.83$ \\
$w5k$ & 0.75 & 62.70 & $+2.86$ \\
$w5k$ & 1.00 & 62.79 & $+2.95$ \\
\midrule
$w0$  & 0.25 & 62.77 & $+2.93$ \\
$w0$  & 0.50 & 62.56 & $+2.72$ \\
$w0$  & 0.75 & 63.07 & $+3.23$ \\
$w0$  & 1.00 & 62.48 & $+2.64$ \\
\bottomrule
\end{tabular}
\end{table}

\paragraph{Three-seed validation.}
Table~\ref{tab:ct-phase2-3seeds} reports per-seed best validation mIoU,
mean, standard deviation, and $\Delta$ versus the 3-seed baseline mean
for all validated configurations.

\begin{table}[H]
\centering
\caption{Phase~II three-seed validation (N\,=\,372, best-observed across
runs). Baseline: PAD-MTL (Table~5, Multi-Task), $\lambda_\text{ct}=0$.
$\Delta = \text{mean}(\%) - 60.75$\,pp.}
\label{tab:ct-phase2-3seeds}
\renewcommand{\arraystretch}{1.3}
\small
\begin{tabularx}{\textwidth}{Xcccccr}
\toprule
Configuration &
Seed~7 & Seed~25 & Seed~42 & Mean & Std & $\Delta$\,(pp) \\
\midrule
Baseline (PAD-MTL) &
59.84 & 60.92 & 61.50 & 60.75 & 0.69 & -- \\
\midrule
single cosine, $w0$, $\lambda=0.75$ &
63.08 & 62.89 & 62.48 & 62.82 & 0.25 & $+2.07$ \\
single cosine, $w0$, $\lambda=1.00$ &
63.01 & 63.37 & 63.37 & 63.25 & 0.17 & $+2.50$ \\
single cosine, $w0$, $\lambda=1.25$ &
62.44 & 62.24 & 64.33 & 63.00 & 0.94 & $+2.25$ \\
\midrule
single MSE, $w5k$, $\lambda=1.00$ &
61.32 & 62.92 & 63.93 & 62.72 & 1.07 & $+1.97$ \\
single MSE, $w0$, $\lambda=1.00$ &
62.06 & 62.91 & 64.91 & 63.29 & 1.20 & $+2.54$ \\
\midrule
dual cosine, $w0$, $\lambda=0.50$ &
62.54 & 62.91 & 63.15 & 62.87 & 0.25 & $+2.12$ \\
dual cosine, $w0$, $\lambda=0.75$ &
63.22 & 62.91 & 62.73 & 62.95 & 0.20 & $+2.20$ \\
dual cosine, $w0$, $\lambda=1.00$ &
62.77 & 61.81 & 63.63 & 62.74 & 0.74 & $+1.99$ \\
dual cosine, $w5k$, $\lambda=0.50$ &
62.13 & 63.15 & 61.58 & 62.29 & 0.65 & $+1.54$ \\
\bottomrule
\end{tabularx}
\end{table}

\noindent
\textbf{Phase-II findings.}
\begin{itemize}
\item All validated configurations improve over the baseline mean
  ($60.75$\,\%), with gains ranging from $+1.54$ to $+2.54$\,pp. This
  confirms that $\mathcal{L}_\text{ct}$ provides consistent positive
  signal across diverse hyperparameter choices.

\item The most stable configuration is single cosine, $w0$,
  $\lambda_\text{ct}=1.00$: mean $63.25$\,\%, $\Delta=+2.50$\,pp, and the
  lowest standard deviation ($0.17$\,\%) across all variants, indicating
  highly reliable improvement over the baseline.

\item Single MSE, $w0$, $\lambda_\text{ct}=1.00$ achieves the highest
  displayed mean ($63.29$\,\%, $\Delta=+2.54$\,pp), but with considerably
  larger variance (std $1.20$\,\%), reflecting higher seed-to-seed
  instability.

\item Dual-projection results consistently fall below the single-projection
  optimum. The best dual variant (dual cosine, $w0$,
  $\lambda_\text{ct}=0.75$) reaches mean $62.95$\,\% ($+2.20$\,pp), approximately
  $0.30$\,pp below the best single configuration.

\item The single-seed Phase~I peak for dual cosine
  ($w0$, $\lambda_\text{ct}=0.50$: $63.66$\,\%) does not generalise to the
  three-seed mean ($62.87$\,\%), underscoring the importance of multi-seed
  validation before concluding on optimal $\lambda_\text{ct}$.

\item Dual cosine with $w5k$ yields the smallest gain among all
  configurations ($+1.54$\,pp), confirming that combining warmup with
  dual projection is suboptimal compared to either component alone.
\end{itemize}

% ============================================================
\subsection{Summary and Conclusions}
\label{subsec:ct-summary}

The two-phase evaluation supports the following conclusions:

\begin{enumerate}
\item \textit{Effectiveness.} $\mathcal{L}_\text{ct}$ provides consistent
  positive gains over the PAD-MTL baseline across a range of
  configurations, confirming that explicit cross-task feature alignment
  is beneficial for label-efficient segmentation.

\item \textit{Single projection is preferable.} Single-projection variants
  achieve higher mean mIoU and lower variance than dual-projection
  counterparts.   The single cosine configuration ($w0$,
  $\lambda_\text{ct}=1.00$) is selected as the primary design choice:
  mean $63.25$\,\% ($+2.50$\,pp over baseline), std $0.17$\,\%.

\item \textit{Warmup is not universally required.} For the cosine loss,
  $w0$ consistently outperforms $w5k$. Warmup remains useful for MSE
  under single projection but offers no benefit once dual projection is
  introduced.

\item \textit{Robust $\lambda_\text{ct}$ range.} For single cosine with
  $w0$, $\lambda_\text{ct}=1.00$ offers the best balance of performance
  and stability across three seeds.
\end{enumerate}
